\subsection{Density Theorem}

THEOREM 3 (Jacobson). --- Let M be a semisimple A-module, and let c be an endomorphism of the additive group M. Then c belongs to the bicommutant \(A_{M}^{\prime\prime}\) of M if and only if it satisfies the following condition:

\begin{enumerate}
\def\labelenumi{(\Alph{enumi})}
\setcounter{enumi}{3}
\tightlist
\item
  For every finite subset \(\mathrm{F}\) of \(\mathrm{M}\) , there exists an element \(a\) of \(\mathrm{A}\) such that \(c\) coincides with \(a_{\mathrm{M}}\) on \(\mathrm{F}\) .
\end{enumerate}

First, suppose that \(c\) satisfies condition (D). Let \(u\) be an element of \(\mathrm{A}_{\mathrm{M}}^{\prime}\) . Let \(x\) be an element of \(\mathrm{M}\) , and apply condition (D) to the subset \(\mathrm{F} = \{x, u(x)\}\) . There exists an element \(a\) of \(\mathrm{A}\) such that \(c(x) = ax\) and \(c(u(x)) = au(x)\) , so that \(u(c(x)) = u(ax) = au(x) = c(u(x))\) . Since \(x\) is arbitrary, we have \(cu = uc\) ; as this holds for all \(u\) , we have \(c \in \mathrm{A}_{\mathrm{M}}^{\prime \prime}\) .

For the converse, we use the following lemma.

Lemma 3. --- Let M be a semisimple A-module. Let B be the bicommutant of the A-module M. Then every A-submodule of M is a B-submodule of M.

Let N be an A-submodule of M. By Corollary 2 of VIII, p.~56, there exists a projector p of the A-module M with image N; it belongs to \(A_{M}^{\prime}\) . Since we have the relation pb = bp for every \(b \in B\) , we obtain that N is a B-submodule of M.

Let us conclude the proof of Theorem 3. Suppose that c belongs to \(A_{M}^{\prime\prime}\) , and let \(F = \{x_{1}, \ldots, x_{n}\}\) be a finite subset of M. Denote the element \((x_{1}, \ldots, x_{n})\) of \(M^{n}\) by x. The A-module \(M^{n}\) is semisimple and, by Proposition 2 of VIII, p.~79, its bicommutant coincides with the homotheties of the \(A_{M}^{\prime\prime}\) -module \(M^{n}\) . By Lemma 3, the A-submodule Ax of \(M^{n}\) is an \(A_{M}^{\prime\prime}\) -submodule of \(M^{n}\) . Let \(a \in A\) be such that \((cx_{1}, \ldots, cx_{n})\) is equal to ax. Then c coincides with \(a_{M}\) on \(\{x_{1}, \ldots, x_{n}\}\) , which implies condition (D).

Remark. --- Denote the endomorphism ring of the additive group of M by E. Endow M with the discrete topology; the ring E consists of mappings from M to M, and we can endow it with the topology induced by the product topology on \(M^{M}\) (``topology of simple convergence in M'', Gen.~Top., I, §2, No.~3, p.~31). The topology on E is Hausdorff and compatible with the additive group structure on E. For every f in E, the mappings \(g \mapsto f \circ g\) and \(g \mapsto g \circ f\) from E to E are continuous. Consequently, the commutant of every subset of E is closed in E. Theorem 3 therefore implies that \(A_{M}^{\prime\prime}\) is the closure of \(A_{M}\) in E.

